% Compile with: uplatex file.tex && dvipdfmx file.dvi
%           or: latexmk -uplatex file.tex
%\documentclass[uplatex,dvipdfmx]{jsarticle}
\documentclass[dvipdfmx,a4paper,11pt]{jsarticle}
\usepackage{amsmath, amssymb}
\usepackage{tikz}

\title{scenario\_coordinate\_jacobian}
\date{\today}

\begin{document}
\maketitle

単位円板 $x^2+y^2\le1$ 上で、被積分関数 $x^2+y^2$ の範囲つき定積分を登録する。この時点ではまだ CAS には渡さず、表示専用の未評価オブジェクトとして保持する。

$x^2$+$y^2$$\le$1 の領域で $x^2$+$y^2$ を定積分

\begin{equation}
  \iint_{x^2+y^2 \le 1} (x^2+y^2) \, dx \, dy
  \label{eq:E1}
\end{equation}

この定積分をデカルト座標 $(x,y)$ から極座標 $(r,\theta)$ へ変数変換する（$x=r\cos\theta,\ y=r\sin\theta$）。ヤコビアン行列式 $|J|=r$ が被積分関数に掛かり、積分範囲も新しい変数のものへ自動的に書き換えられる。

E1 を x,y から r,theta へ x=r*cos(theta), y=r*sin(theta) で変数変換して

\begin{equation}
  \int_{0}^{2\pi} \int_{0}^{1} r^3 \, dr \, d\theta
  \label{eq:E2}
\end{equation}

変数変換後の定積分を実際に評価し、数値を得る。

この定積分を評価して

\begin{equation}
  \int_{0}^{2\pi} \int_{0}^{1} r^3 \, dr \, d\theta = \frac{\pi}{2}
  \label{eq:E3}
\end{equation}

今度は単位球 $x^2+y^2+z^2\le1$ 上で、被積分関数を $1$（体積そのもの）とする範囲つき定積分を、表示専用オブジェクトとして登録する。

$x^2$+$y^2$+$z^2$$\le$1 の領域で 1 を定積分

\begin{equation}
  \iiint_{x^2+y^2+z^2 \le 1} (1) \, dx \, dy \, dz
  \label{eq:E4}
\end{equation}

デカルト座標から球座標 $(r,\theta,\phi)$ へ変数変換する（$x=r\sin\theta\cos\phi,\ y=r\sin\theta\sin\phi,\ z=r\cos\theta$）。ヤコビアン行列式 $|J|=r^2\sin\theta$ が体積要素に現れる。

E4 を x,y,z から r,theta,phi へ x=r*sin(theta)*cos(phi), y=r*sin(theta)*sin(phi), z=r*cos(theta) で変数変換して

\begin{equation}
  \int_{0}^{2\pi} \int_{0}^{\pi} \int_{0}^{1} r^2 \sin \vartheta \, dr \, d\theta \, d\phi
  \label{eq:E5}
\end{equation}

変数変換後の定積分を評価し、単位球の体積を得る。

この定積分を評価して

\begin{equation}
  \int_{0}^{2\pi} \int_{0}^{\pi} \int_{0}^{1} r^2 \sin \vartheta \, dr \, d\theta \, d\phi = \frac{4 \pi}{3}
  \label{eq:E6}
\end{equation}

今度は座標変換の公式が定型でない一般の例として、$\int_0^1\int_0^1(x+y)\,dx\,dy$ を LaTeX 記法でそのまま入力し、表示専用オブジェクトとして登録する。

$\int_{0}^{1} \int_{0}^{1} (x+y) \, dx \, dy$

\begin{equation}
  \int_{0}^{1} \int_{0}^{1} x+y \, dx \, dy
  \label{eq:E7}
\end{equation}

$(x,y)$ から $(u,v)$ への一次変換 $x=2u+v,\ y=u+3v$（ヤコビアン行列式 $\det J=5$）で変数変換する。標準的な座標変換ではないため、新しい積分範囲 $u:0..1,\ v:0..1$ をここで利用者が明示的に指定する。

E7 を x,y から u:0..1, v:0..1 へ x=2*u+v, y=u+3*v で変数変換して

\begin{equation}
  \int_{0}^{1} \int_{0}^{1} 20 v+15 u \, dv \, du
  \label{eq:E8}
\end{equation}

変数変換後の定積分を評価し、数値を得る。

この定積分を評価して

\begin{equation}
  \int_{0}^{1} \int_{0}^{1} 20 v+15 u \, dv \, du = \frac{35}{2}
  \label{eq:E9}
\end{equation}

\end{document}